% Section 10.1, from lectures/L10/appendix/a_device_model.md.

\section{The Device Model}
\label{c10:app:a}

How a driver finds its hardware without being told where it is, and what the kernel does with the
three-way relationship between buses, devices and drivers.

\S\ref{c10:app:b} is where the description comes from.

The idea to carry: \textbf{you do not write code that goes looking for hardware.} You register a
driver that describes what it can handle, something else registers a device that describes what
exists, and the kernel calls your \code{probe} when the two match. Everything in this section is a
consequence of that inversion.

\subsection{Three things and a match}
\label{c10:app:a1}

\begin{booktable}{@{}l>{\hsize=1.28\hsize}L>{\hsize=0.72\hsize}L@{}}
   & \thead{Is} & \thead{Registered by} \\
  \midrule
  \textbf{Bus} & A matching policy & The kernel, once \\
  \textbf{Device} & Something that exists & Whatever discovered it \\
  \textbf{Driver} & Code that can handle a class of devices & Your module \\
\end{booktable}

\noindent Each bus keeps two lists, its devices and its drivers. Whenever either list changes, the
bus tries to match the new entry against every entry on the other list, and calls \code{probe} on
each success.

\textbf{A bus is a matching policy, not a wire.} \code{platform} is a bus and there is no platform
bus in any schematic; it is the bus for devices that cannot be discovered, which is most of an SoC.
PCI is a bus and can enumerate itself. USB is a bus and enumerates on hotplug. What they share is
the matching machinery, not the electricals.

The consequence people find surprising is the ordering: \textbf{the device and the driver can arrive
in either order.} On the target, the device for \code{qa-dev} exists from boot and sits there
unbound:

\begin{console}
# ls /sys/bus/platform/devices/c000000.qa-dev
driver_override  modalias  of_node  power  subsystem  uevent
# ls /sys/bus/platform/devices/c000000.qa-dev/driver
ls: no such file
\end{console}

\noindent Loading a module that registers a matching driver binds it immediately. Load the module
first on a machine where the device appears later, and it binds then instead. Your \code{probe} is
written the same way either way.

\lecturefigure{lectures/L10/appendix/images/device_model.png}
  {A device from the device tree and a driver with its match table, both pointing down at the platform
   bus, which compares compatible against every registered driver. A red arrow marked match leads to
   probe, with a note that everything probe takes with devm\_ is released when the bind ends.}
  {fig:device_model}

\subsection{The tree, in sysfs}
\label{c10:app:a2}

\file{/sys} is the device model made visible, which is why \S\ref{c2:app:b4} could describe it as
generated rather than written.

\begin{narrowtable}{@{}ll@{}}
  \thead{Path} & \thead{Holds} \\
  \midrule
  \code{/sys/devices/} & The real tree, by physical topology \\
  \code{/sys/bus/*/devices/} & Symlinks, grouped by bus \\
  \code{/sys/bus/*/drivers/} & One directory per registered driver \\
  \code{/sys/class/*/} & Symlinks, grouped by \textbf{what a device is} \\
\end{narrowtable}

\noindent The distinction between the last two is worth having. \file{/sys/bus} groups by \emph{how
it is attached}; \file{/sys/class} groups by \emph{what it does}. A serial port reached over USB
appears under \file{/sys/bus/usb} and under \file{/sys/class/tty}, and a program looking for serial
ports wants the second. Chapter~\ref{c11:ch} is about getting into \file{/sys/class} at all.

Every device and driver directory there is a \textbf{kobject}, and every kobject is reference
counted. That is what makes it safe for a process to hold a device open while somebody unbinds its
driver: the memory is not freed until the last reference goes.

\subsection{\code{probe} and \code{remove}}
\label{c10:app:a3}

\begin{ccode}
static int qa_probe(struct platform_device *pdev);
static void qa_remove(struct platform_device *pdev);
\end{ccode}

\noindent \code{probe} is called once per matching device, and \textbf{it can be called more than
once}: two of the same chip on a board means two calls. So everything it allocates must hang off the
device it was given, not off a static in your module. A driver with a global \code{regs} pointer
works perfectly until the second device appears.

That is what \code{platform_set_drvdata} and \code{dev_get_drvdata} are for: one private structure
per device, retrieved by any callback that is handed the \code{struct device}.

\code{remove} must undo what \code{probe} did, and returns \code{void} in 6.12: \textbf{it is not
allowed to fail.} By the time it is called the decision has been made, and there is nothing sensible
for the driver core to do with an error.

\subsubsection{What may be done in probe}
\label{c10:app:a:what-may-be-done-in-probe}

\code{probe} runs in process context and may sleep, which means \code{GFP_KERNEL}, mutexes and I/O
are all available. Two things to know:
\begin{itemize}
  \item \textbf{Do not assume a fixed order between devices.} If your device needs a clock provided
    by another driver that has not probed yet, return \code{-EPROBE_DEFER} and the core will call
    you again later. Returning a real error instead makes your driver depend on link order, which is
    not a thing you control.
  \item \textbf{Publish last.} Anything that makes the device reachable from userspace, a character
    device or a sysfs attribute, should be the last thing \code{probe} does, for the same reason
    \S\ref{c5:app:a3} gave about \code{cdev_add}.
\end{itemize}

\subsection{\code{devm_}, and what it is really attached to}
\label{c10:app:a4}

\S\ref{c6:app:a7} introduced these and could not use them, because a bare module has no \code{struct
device}. \code{probe} is handed one, so this chapter can.

\begin{ccode}
priv       = devm_kzalloc(&pdev->dev, sizeof(*priv), GFP_KERNEL);
priv->regs = devm_platform_ioremap_resource(pdev, 0);
ret        = devm_request_irq(&pdev->dev, irq, qa_isr, 0, dev_name(&pdev->dev), priv);
\end{ccode}

\noindent Three acquisitions, no error labels, and no \code{remove} code for any of them. Compare
\S\ref{c6:app:b7}, which is the same work with a \code{goto} ladder underneath it.

\textbf{The release happens when the device is unbound from the driver.} Not at module unload, and
not at process exit. That distinction is testable, and the lab tests it:

\begin{console}
# echo c000000.qa-dev > /sys/bus/platform/drivers/qa_platform/unbind
qa_platform c000000.qa-dev: removed after 1296 interrupts, 1297 samples

# grep -c c000000.qa-dev /proc/iomem
0
# lsmod | grep qa_platform
qa_platform 12288 0 - Live
\end{console}

\noindent The module is \textbf{still loaded}. The region and the interrupt are gone, released by
the driver core walking the device's list of managed resources in reverse, and so are the sysfs
attributes, which the core removes itself because they came from the driver's \code{.dev_groups}.
Bind it again and \code{probe} runs again and they come back.

That is the whole mechanism: \code{devm_} builds a list of destructors attached to one device's
binding. Knowing that is what stops the two surprises, which are things released earlier than
expected (because something unbound) and things not released at all (because they were attached to
the wrong device, or acquired without \code{devm_} at all).

\subsection{Matching, and how a module gets loaded}
\label{c10:app:a5}

\begin{ccode}
static const struct of_device_id qa_of_match[] = {
        { .compatible = "qacademy,qa-dev-1.0" },
        { }
};
MODULE_DEVICE_TABLE(of, qa_of_match);
\end{ccode}

\noindent The table is the driver's statement of what it handles, terminated by an empty entry. The
\code{MODULE_DEVICE_TABLE} line does something different and is easy to omit: it copies those
strings into the module's binary metadata, so that they can be found \textbf{without loading the
module}.

Four steps turn a device appearing into a module loading, and none of them is magic:
\begin{enumerate}
  \item \code{MODULE_DEVICE_TABLE} puts the compatible strings in the \code{.ko}'s \code{.modinfo}
    section.
  \item \code{depmod} reads every module's metadata and writes
    \code{/lib/modules/<version>/modules.alias}.
  \item The device's \code{uevent} carries a \code{MODALIAS=} string built from its own
    \code{compatible}.
  \item udev matches one against the other and runs \code{modprobe}.
\end{enumerate}
You can see step 3 on the target:

\begin{console}
# cat /sys/bus/platform/devices/c000000.qa-dev/uevent
DRIVER=qa_platform
OF_NAME=qa-dev
OF_FULLNAME=/platform-bus@c000000/qa-dev@0
OF_COMPATIBLE_0=qacademy,qa-dev-1.0
OF_COMPATIBLE_N=1
MODALIAS=of:Nqa-devT(null)Cqacademy,qa-dev-1.0
\end{console}

\noindent This target has no udev, so nothing performs step 4 and modules are loaded by hand. That
does not change steps 1 to 3, and it is worth knowing which part is missing rather than concluding
that autoloading is mysterious.

\subsection{Whose name is it}
\label{c10:app:a6}

A detail that makes the inversion concrete. In Chapter~\ref{c6:ch} the driver chose the name it
registered its region under:

\begin{console}
0c000000-0c000fff : qa_mmio
\end{console}

\noindent As a platform driver it does not choose, and the entry reads:

\begin{console}
0c000000-0c000fff : c000000.qa-dev qa-dev@0
 17:        243          0  GIC-0 144 Level     c000000.qa-dev
\end{console}

\noindent The device is named from its \textbf{translated address} and its node name, by the driver
core, from the device tree. Two consequences: the same driver bound to two instances produces two
distinguishable entries, which a hand-chosen string could not; and \code{dev_name(&pdev->dev)} is
the right thing to pass wherever a name is wanted, including to \code{devm_request_irq}.

The same applies to messages. \code{dev_info(&pdev->dev, ...)} prints \code{qa_platform
c000000.qa-dev: probed, irq 17} and identifies which device spoke; Chapter~\ref{c4:ch}'s
\code{pr_info} cannot, and that stops being a stylistic preference the moment there are two of
anything.

\subsection{Binding by hand}
\label{c10:app:a7}

Every driver directory has \code{bind} and \code{unbind}:

\begin{console}
# echo c000000.qa-dev > /sys/bus/platform/drivers/qa_platform/unbind
# echo c000000.qa-dev > /sys/bus/platform/drivers/qa_platform/bind
\end{console}

\noindent This is genuinely useful and not just a demonstration. It is how you hand a device to
\code{vfio-platform} for passthrough, how you reset a wedged driver without rebooting, and how you
test that your \code{remove} path actually works, which is otherwise the least-exercised code in a
driver.

\code{driver_override} in the device's directory forces a specific driver regardless of matching,
which is the other half of the passthrough workflow.

\textbf{Unbind while userspace holds the device open} is the case that catches people. The kobject
reference count keeps the memory alive, but your \code{remove} has already run, so any pointer the
file operations still hold is pointing at released hardware state. Handling that properly is what a
driver's reference counting is for, and it is why frameworks (Chapter~\ref{c11:ch}) are worth using:
they have already solved it.
